\documentclass{article}
\usepackage[T1]{fontenc}
\usepackage{lmodern}
\usepackage{graphicx}
\usepackage{amsmath,amssymb}
\usepackage[a4paper,margin=2cm]{geometry}
\usepackage[absolute,overlay]{textpos}
\setlength{\TPHorizModule}{1mm}
\setlength{\TPVertModule}{1mm}
\usepackage[indonesian]{babel}
\usepackage{hyperref}
\setlength{\emergencystretch}{2em}
% unit: Halaman 36

\begin{document}

% === Halaman 36 (python/native) ===

36

Kriteria Steganografi yang Bagus

1. Imperceptible 


Keberadaan pesan rahasia tidak dapat dipersepsi secara visual atau secara 

audial

2.   Fidelity.


 Kualitas cover-object  tidak jauh berubah akibat penyisipan pesan rahasia. 

3.    Recovery. 


 Pesan yang disembunyikan harus dapat diekstraksi kembali. 

4.    Capacity 


 Ukuran pesan yang disembunyikan sedapat mungkin besar

Catatan: Robustnes bukan isu penting di dalam steganografi

\end{document}
